\documentclass[11pt]{article}
\usepackage{cornell-notes}
\usepackage{mathtools}
\usepackage{needspace}
\usepackage{ragged2e}
\tcbuselibrary{most}
\geometry{letterpaper,margin=0.72in,headheight=15pt}
\setCornellCollection{Cloud Architecture Cornell Notes}
\setCornellUnitType{Chapter}
\setCornellUnitNumber{13}
\setCornellUnitTitle{Fundamental Cloud Architectures}

\definecolor{Navy}{HTML}{17324D}
\definecolor{Teal}{HTML}{147D78}
\definecolor{Sky}{HTML}{E8F2F7}
\definecolor{CueGray}{HTML}{EEF2F4}
\definecolor{Warm}{HTML}{FFF3D6}
\definecolor{Risk}{HTML}{953B35}
\definecolor{Green}{HTML}{276749}
\definecolor{RuleGray}{HTML}{AAB8C2}
\definecolor{Text}{HTML}{25313A}

\hypersetup{colorlinks=true,linkcolor=Navy,urlcolor=Teal,citecolor=Navy,pdfborder={0 0 0}}
\color{Text}
\setlength{\parindent}{0pt}
\setlength{\parskip}{4pt}
\setlength{\emergencystretch}{3em}
\setlength{\LTpre}{3pt}
\setlength{\LTpost}{5pt}
\renewcommand{\arraystretch}{1.16}
\setlist{nosep,leftmargin=1.45em}
\hfuzz=0.2pt
\hbadness=2500

\titleformat{\section}{\Large\bfseries\color{Navy}\RaggedRight}{}{0pt}{}[\color{Teal}\titlerule]
\titleformat{\subsection}{\large\bfseries\color{Teal}\RaggedRight}{}{0pt}{}
\titlespacing*{\section}{0pt}{1.35ex plus .4ex}{.7ex}
\titlespacing*{\subsection}{0pt}{1.1ex plus .3ex}{.45ex}

\pagestyle{fancy}
\fancyhf{}
\fancyhead[L]{\small\color{Navy}\textbf{Cloud Architecture Cornell Notes}}
\fancyhead[R]{\small\color{Teal}\ChapterMark}
\fancyfoot[L]{\footnotesize\color{Navy}Concepts $\bullet$ Mechanisms $\bullet$ Evidence $\bullet$ Decisions}
\fancyfoot[R]{\footnotesize\color{Navy}\thepage\ of \pageref*{LastPage}}
\renewcommand{\headrulewidth}{0.5pt}
\renewcommand{\footrulewidth}{0.35pt}

\newtcolorbox{orientationbox}{enhanced,breakable,colback=Sky,colframe=Teal,boxrule=.75pt,arc=2mm,title={Chapter Orientation},fonttitle=\bfseries,left=2mm,right=2mm,top=1.5mm,bottom=1.5mm}
\newtcolorbox{sectionsummarybox}{enhanced,breakable,colback=CueGray,colframe=RuleGray,boxrule=.55pt,arc=1.5mm,title={Topic Summary},fonttitle=\bfseries\color{Navy},left=2mm,right=2mm,top=1mm,bottom=1mm}
\newtcolorbox{practicebox}{enhanced,breakable,colback=white,colframe=Green,boxrule=.75pt,arc=2mm,title={Practitioner Lens},fonttitle=\bfseries}
\newtcolorbox{scopenotebox}{enhanced,breakable,colback=Warm,colframe=Risk,boxrule=.7pt,arc=2mm,title={Terminology and Scope Note},fonttitle=\bfseries}
\newtcolorbox{synthesisbox}{enhanced,breakable,colback=Navy!5,colframe=Navy,boxrule=.85pt,arc=2mm,title={Chapter Synthesis},fonttitle=\bfseries}
\newtcolorbox{takeawaybox}{enhanced,breakable,colback=Teal!8,colframe=Teal,boxrule=.9pt,arc=2mm,title={One-Sentence Takeaway},fonttitle=\bfseries}
\newcommand{\CornellHeader}{\rowcolor{Navy}\color{white}\bfseries Cue / Active Recall & \color{white}\bfseries Notes / Reasoning \\\endhead}
\newcommand{\CornellRow}[2]{\cellcolor{CueGray}\textbf{#1} & #2 \\\addlinespace[3pt]}

\newcommand{\ChapterMark}{Chapter 13}
\title{Chapter 13: Fundamental Cloud Architectures\\[2mm]\large Cornell Notes}
\author{}
\date{}
\hypersetup{pdftitle={Chapter 13: Fundamental Cloud Architectures - Cornell Notes},pdfauthor={},pdfsubject={Cloud computing and architecture study notes},pdfkeywords={cloud, architecture, security, operations, Cornell notes}}
\begin{document}
\maketitle
\thispagestyle{fancy}
\vspace{-1.5em}
\begin{orientationbox}
\textbf{Central problem.} Explain foundational architecture patterns for workload distribution, pooling, dynamic and elastic scaling, service balancing, bursting, storage provisioning, storage redundancy, and multicloud use.\par\smallskip
\textbf{Prerequisites.} Infrastructure and specialized mechanisms, load balancing, resource pools, monitoring, replication, virtualization, storage, networking, and deployment models.\par\smallskip
\textbf{Connection.} These architectures are the first composition layer: they arrange individual mechanisms into repeatable solutions that advanced and specialized patterns extend.
\end{orientationbox}
\section{Learning Objectives}
\begin{enumerate}
\item Define the principal concepts and decision boundaries in Fundamental Cloud Architectures.
\item Explain how 13.1 Workload distribution architecture contributes to the chapter's central problem.
\item Distinguish Workload distribution from Resource pool using responsibility, behavior, and evidence.
\item Trace components, data, control, identity, capacity, and failure across 13.2 Resource pooling architecture.
\item Analyze benefits together with assumptions, limitations, security risks, and operational tradeoffs.
\item Evaluate provider and consumer responsibilities at every relevant trust or administrative boundary.
\item Select verification evidence for architecture, configuration, monitoring, cost, and service-quality claims.
\item Apply the chapter's reasoning to a realistic design, review, migration, or incident scenario.
\end{enumerate}
\section{Cornell Cue-and-Notes}
\Needspace{9\baselineskip}
\subsection{13.1 Workload distribution architecture}
\begin{longtable}{@{}C N@{}}
\CornellHeader
\CornellRow{What should be remembered about Workload distribution?}{A load balancer receives requests and selects among available resource instances using health, weights, observed load, affinity, or another policy. \textbf{Review lens.} Model session state, health checks, uneven work, saturation, instance removal, balancer failure, and capacity after losing a member.}
\end{longtable}
\begin{sectionsummarybox}
Identical resource instances and a load balancer horizontally distribute workload to reduce overuse, underuse, and single-instance pressure.
\end{sectionsummarybox}
\Needspace{9\baselineskip}
\subsection{13.2 Resource pooling architecture}
\begin{longtable}{@{}C N@{}}
\CornellHeader
\CornellRow{What should be remembered about Resource pool?}{Physical servers, virtual servers, storage, networks, CPUs, memory, or other resources can be synchronized and allocated through a resource-management system. \textbf{Review lens.} Define membership, capacity, tenancy, allocation, release, synchronization, quotas, placement, health, and pool exhaustion.}
\end{longtable}
\begin{sectionsummarybox}
Similar resources are grouped into a managed pool from which consumers or automation receive capacity as needed.
\end{sectionsummarybox}
\Needspace{9\baselineskip}
\subsection{13.3 Dynamic scalability architecture}
\begin{longtable}{@{}C N@{}}
\CornellHeader
\CornellRow{What should be remembered about Dynamic scalability?}{A usage monitor supplies measurements, an automated scaling listener evaluates thresholds, and resource management allocates or reclaims instances while a load balancer adjusts routing. \textbf{Review lens.} Test metric delay, cooldowns, provisioning time, state, quotas, minimum capacity, maximum scale, oscillation, scale-in safety, and cost.}
\end{longtable}
\begin{sectionsummarybox}
Predefined workload conditions trigger automatic addition or removal of resources from pools so capacity follows demand.
\end{sectionsummarybox}
\Needspace{9\baselineskip}
\subsection{13.4 Elastic resource capacity architecture}
\begin{longtable}{@{}C N@{}}
\CornellHeader
\CornellRow{What should be remembered about Elastic resource capacity?}{Management and automation dynamically resize vertical capacity using resource pools, allowing a workload to scale without necessarily adding a separate service instance. \textbf{Review lens.} Identify hot-resize support, disruption, guest recognition, maximum host capacity, contention, rollback, billing granularity, and the remaining single-instance risk.}
\end{longtable}
\begin{sectionsummarybox}
An existing virtual server gains or releases CPU and memory in response to changing processing requirements.
\end{sectionsummarybox}
\Needspace{9\baselineskip}
\subsection{13.5 Service load balancing architecture}
\begin{longtable}{@{}C N@{}}
\CornellHeader
\CornellRow{What should be remembered about Service load balancing?}{The architecture specializes workload distribution for service endpoints and may use external or built-in balancing with a pool of equivalent service implementations. \textbf{Review lens.} Verify endpoint identity, version compatibility, statelessness or shared state, health semantics, retries, idempotency, and balanced dependency load.}
\end{longtable}
\begin{sectionsummarybox}
Redundant cloud-service implementations are placed behind load-balancing logic so service requests can be distributed across hosts or instances.
\end{sectionsummarybox}
\Needspace{9\baselineskip}
\subsection{13.6 Cloud bursting architecture}
\begin{longtable}{@{}C N@{}}
\CornellHeader
\CornellRow{What should be remembered about Cloud bursting?}{Monitoring detects local saturation, cloud capacity becomes active, traffic or work is redirected, and remote resources are released after the burst. \textbf{Review lens.} Assess network latency and capacity, data placement, security equivalence, identity, licensing, state, warm-up time, provider quotas, cost, and burst-in safety.}
\end{longtable}
\begin{sectionsummarybox}
On-premises capacity expands into prearranged cloud resources when local thresholds are reached and contracts when demand falls.
\end{sectionsummarybox}
\Needspace{9\baselineskip}
\subsection{13.7 Elastic disk provisioning architecture}
\begin{longtable}{@{}C N@{}}
\CornellHeader
\CornellRow{What should be remembered about Elastic disk provisioning?}{Thin or elastic provisioning exposes an apparent capacity while physical blocks are allocated on demand and usage is monitored. \textbf{Review lens.} Monitor pool exhaustion, overcommit, growth limits, reclaim behavior, performance, snapshots, billing basis, alerts, and emergency expansion.}
\end{longtable}
\begin{sectionsummarybox}
Logical storage presented to a consumer grows from a shared pool as actual data consumption increases, reducing waste from fixed preallocation.
\end{sectionsummarybox}
\Needspace{9\baselineskip}
\subsection{13.8 Redundant storage architecture}
\begin{longtable}{@{}C N@{}}
\CornellHeader
\CornellRow{What should be remembered about Redundant storage?}{A gateway or storage system coordinates copies and redirects access when a primary path or device fails, subject to consistency and recovery policy. \textbf{Review lens.} Define synchronous or asynchronous behavior, failure-domain separation, lag, promotion, split brain, corruption propagation, encryption, and restoration tests.}
\end{longtable}
\begin{sectionsummarybox}
Data is replicated across storage resources so a device, controller, network, site, or security event does not cause unacceptable loss or interruption.
\end{sectionsummarybox}
\Needspace{9\baselineskip}
\subsection{13.9 Multicloud architecture}
\begin{longtable}{@{}C N@{}}
\CornellHeader
\CornellRow{What should be remembered about Multicloud?}{Workloads, data, or supporting services are allocated across providers whose delivery models, interfaces, controls, regions, and operating practices can differ. \textbf{Review lens.} Evaluate portability, identity, networking, data movement, observability, policy consistency, skills, SLA composition, cost, exit, and added complexity.}
\end{longtable}
\begin{sectionsummarybox}
Services or resources from two or more public clouds are used within one consumer strategy to obtain capability, geography, negotiation, or dependency benefits.
\end{sectionsummarybox}
\Needspace{9\baselineskip}
\subsection{13.10 Case application}
\begin{longtable}{@{}C N@{}}
\CornellHeader
\CornellRow{What should be remembered about Architecture selection?}{The practitioner identifies the problem, constraints, candidate pattern, participating mechanisms, expected improvement, tradeoffs, and operating evidence. \textbf{Review lens.} Use a decision record and validate the highest-risk assumption with load, failure, recovery, and cost tests.}
\end{longtable}
\begin{sectionsummarybox}
A case application selects an architecture only after relating workload and business conditions to mechanisms, quality attributes, cost, and risk.
\end{sectionsummarybox}
\begin{practicebox}
For each model, mechanism, architecture, metric, or case decision, record the objective, owner, trust boundary, inputs, dependencies, expected behavior, failure consequence, measurement point, response, and retained evidence. A diagram or product label is a hypothesis until runtime and recovery tests support it.
\end{practicebox}
\section{Chapter Recap Matrix}
\begingroup\small\setlength{\tabcolsep}{4pt}
\begin{longtable}{@{}>{\RaggedRight\arraybackslash}p{0.235\textwidth}>{\RaggedRight\arraybackslash}p{0.37\textwidth}>{\RaggedRight\arraybackslash}p{0.325\textwidth}@{}}
\toprule\rowcolor{Navy}\color{white}\textbf{Element} & \color{white}\textbf{Purpose or Behavior} & \color{white}\textbf{Decision and Evidence}\\\midrule\endfirsthead
\toprule\rowcolor{Navy}\color{white}\textbf{Element} & \color{white}\textbf{Purpose or Behavior} & \color{white}\textbf{Decision and Evidence}\\\midrule\endhead
Workload distribution & A load balancer receives requests and selects among available resource instances using health, weights, observed load, affinity, or another policy. & Model session state, health checks, uneven work, saturation, instance removal, balancer failure, and capacity after losing a member. \\\midrule
Resource pool & Physical servers, virtual servers, storage, networks, CPUs, memory, or other resources can be synchronized and allocated through a resource-management system. & Define membership, capacity, tenancy, allocation, release, synchronization, quotas, placement, health, and pool exhaustion. \\\midrule
Dynamic scalability & A usage monitor supplies measurements, an automated scaling listener evaluates thresholds, and resource management allocates or reclaims instances while a load balancer adjusts routing. & Test metric delay, cooldowns, provisioning time, state, quotas, minimum capacity, maximum scale, oscillation, scale-in safety, and cost. \\\midrule
Elastic resource capacity & Management and automation dynamically resize vertical capacity using resource pools, allowing a workload to scale without necessarily adding a separate service instance. & Identify hot-resize support, disruption, guest recognition, maximum host capacity, contention, rollback, billing granularity, and the remaining single-instance risk. \\\midrule
Service load balancing & The architecture specializes workload distribution for service endpoints and may use external or built-in balancing with a pool of equivalent service implementations. & Verify endpoint identity, version compatibility, statelessness or shared state, health semantics, retries, idempotency, and balanced dependency load. \\\midrule
Cloud bursting & Monitoring detects local saturation, cloud capacity becomes active, traffic or work is redirected, and remote resources are released after the burst. & Assess network latency and capacity, data placement, security equivalence, identity, licensing, state, warm-up time, provider quotas, cost, and burst-in safety. \\\midrule
Elastic disk provisioning & Thin or elastic provisioning exposes an apparent capacity while physical blocks are allocated on demand and usage is monitored. & Monitor pool exhaustion, overcommit, growth limits, reclaim behavior, performance, snapshots, billing basis, alerts, and emergency expansion. \\\midrule
Redundant storage & A gateway or storage system coordinates copies and redirects access when a primary path or device fails, subject to consistency and recovery policy. & Define synchronous or asynchronous behavior, failure-domain separation, lag, promotion, split brain, corruption propagation, encryption, and restoration tests. \\\midrule
Multicloud & Workloads, data, or supporting services are allocated across providers whose delivery models, interfaces, controls, regions, and operating practices can differ. & Evaluate portability, identity, networking, data movement, observability, policy consistency, skills, SLA composition, cost, exit, and added complexity. \\\midrule
Architecture selection & The practitioner identifies the problem, constraints, candidate pattern, participating mechanisms, expected improvement, tradeoffs, and operating evidence. & Use a decision record and validate the highest-risk assumption with load, failure, recovery, and cost tests. \\\midrule
\bottomrule
\end{longtable}
\endgroup
\section{Practical Review Checklist}
\begin{longtable}{@{}>{\bfseries\color{Teal}}p{0.16\textwidth}>{\RaggedRight\arraybackslash}p{0.77\textwidth}@{}}
\toprule\textbf{Area} & \textbf{Verification action}\\\midrule\endhead
Scope & $\square$ Define the in-scope services, workloads, users, data, environments, regions, and dependencies for Fundamental Cloud Architectures. \\\addlinespace[2pt]
Roles & $\square$ Name the provider, consumer, owner, administrator, broker, support, security, finance, and risk responsibilities that apply. \\\addlinespace[2pt]
Boundaries & $\square$ Map organizational, trust, network, tenancy, management, data-location, and failure-domain boundaries. \\\addlinespace[2pt]
Architecture & $\square$ Document how Workload distribution, Resource pool, and supporting mechanisms interact through normal, peak, maintenance, and failed states. \\\addlinespace[2pt]
Identity & $\square$ Verify human and workload identities, privileges, credentials, delegation, administrative access, revocation, and audit evidence. \\\addlinespace[2pt]
Data & $\square$ Trace data collection, transfer, processing, storage, replication, backup, retention, export, and deletion. \\\addlinespace[2pt]
Quality & $\square$ Set measurable capacity, latency, availability, reliability, scalability, and recovery criteria relevant to Dynamic scalability. \\\addlinespace[2pt]
Security & $\square$ Challenge the trust and failure assumptions surrounding Elastic resource capacity with abuse, bypass, isolation, and compromise scenarios. \\\addlinespace[2pt]
Operations & $\square$ Confirm monitoring, alert ownership, change control, patching, drift management, incident response, and tested recovery. \\\addlinespace[2pt]
Economics & $\square$ Tie resource behavior to meters, units, rates, budgets, anomalies, full-lifecycle cost, and accountable owners. \\\addlinespace[2pt]
Contracts & $\square$ Reconcile architecture behavior with service scope, formulas, exclusions, evidence, notification, remedies, audit rights, and exit terms. \\\addlinespace[2pt]
Verification & $\square$ Preserve reproducible evidence that Service load balancing and the overall design satisfy their intended outcomes under realistic load and failure. \\\addlinespace[2pt]
\bottomrule
\end{longtable}
\begin{synthesisbox}
Fundamental cloud architectures organize pooled and replicated resources around measurable demand and failure conditions. Distribution, scaling, bursting, storage elasticity, redundancy, and multicloud can improve quality attributes, but each introduces state, synchronization, monitoring, routing, consistency, capacity, and governance decisions that must be made explicit.
\par\smallskip
\textbf{Cross-chapter application.} These architectures are the first composition layer: they arrange individual mechanisms into repeatable solutions that advanced and specialized patterns extend.
\end{synthesisbox}
\section{Self-Test Questions}
\begin{enumerate}
\item What central problem does Fundamental Cloud Architectures address?
\item Define Workload distribution and explain why it matters.
\item Compare Resource pool with Dynamic scalability. Which assumptions and responsibilities differ?
\item Trace the normal behavior or decision path associated with Elastic resource capacity.
\item Which evidence would you request when reviewing Service load balancing?
\item How does this chapter connect to the broader cloud architecture and operating model?
\item Scenario: A team claims that Resource pool is implemented but cannot show current configuration, monitoring, or failure-test evidence. What should the reviewer conclude and do next?
\item Scenario: A workload using Dynamic scalability crosses a provider-consumer boundary and then fails during peak demand. Which responsibilities and failure conditions should be investigated?
\item Why is a named mechanism, model, or architecture insufficient evidence that its intended quality or security outcome is achieved?
\item How would you verify that a change involving Cloud bursting improved the intended outcome without shifting risk or cost elsewhere?
\end{enumerate}
\clearpage
\section{Answer Key}
\begin{enumerate}
\item Explain foundational architecture patterns for workload distribution, pooling, dynamic and elastic scaling, service balancing, bursting, storage provisioning, storage redundancy, and multicloud use.
\item A load balancer receives requests and selects among available resource instances using health, weights, observed load, affinity, or another policy. The review implication is: Model session state, health checks, uneven work, saturation, instance removal, balancer failure, and capacity after losing a member.
\item Physical servers, virtual servers, storage, networks, CPUs, memory, or other resources can be synchronized and allocated through a resource-management system. In contrast, A usage monitor supplies measurements, an automated scaling listener evaluates thresholds, and resource management allocates or reclaims instances while a load balancer adjusts routing. Compare control, ownership, state, dependencies, operating evidence, and failure impact rather than names alone.
\item Management and automation dynamically resize vertical capacity using resource pools, allowing a workload to scale without necessarily adding a separate service instance. A trace should name initiating actor, input or signal, decision point, participating components, state change, output, monitoring, and failure handling.
\item Useful evidence should demonstrate the review lens: Verify endpoint identity, version compatibility, statelessness or shared state, health semantics, retries, idempotency, and balanced dependency load. Include architecture, configuration, event, test, and ownership evidence where applicable.
\item These architectures are the first composition layer: they arrange individual mechanisms into repeatable solutions that advanced and specialized patterns extend.
\item The implementation claim is unsubstantiated. Request the evidence needed to establish current state, then test positive, negative, overload, and failure paths. For this mechanism specifically: Define membership, capacity, tenancy, allocation, release, synchronization, quotas, placement, health, and pool exhaustion.
\item Investigate the agreed responsibility split, boundary controls, capacity and quota, workload state, dependent services, provider telemetry, consumer telemetry, escalation, recovery, and whether the contract measures the consumer-visible failure.
\item A label expresses design intent, not runtime proof. Effectiveness depends on placement, configuration, inputs, state, control independence, capacity, monitoring, response, maintenance, and tested behavior during failure.
\item Establish a before-and-after measure, preserve comparable test conditions, verify the expected behavior, inspect adjacent dependencies and costs, and exercise rollback or recovery. Apply this specific lens: Assess network latency and capacity, data placement, security equivalence, identity, licensing, state, warm-up time, provider quotas, cost, and burst-in safety.
\end{enumerate}
\section{Key-Term Recap}
\begin{longtable}{@{}>{\RaggedRight\arraybackslash}p{0.27\textwidth}>{\RaggedRight\arraybackslash}p{0.66\textwidth}@{}}
\toprule\rowcolor{Navy}\color{white}\textbf{Term} & \color{white}\textbf{Working meaning}\\\midrule\endhead
Workload distribution & A load balancer receives requests and selects among available resource instances using health, weights, observed load, affinity, or another policy. \\\addlinespace[2pt]
Resource pool & Physical servers, virtual servers, storage, networks, CPUs, memory, or other resources can be synchronized and allocated through a resource-management system. \\\addlinespace[2pt]
Dynamic scalability & A usage monitor supplies measurements, an automated scaling listener evaluates thresholds, and resource management allocates or reclaims instances while a load balancer adjusts routing. \\\addlinespace[2pt]
Elastic resource capacity & Management and automation dynamically resize vertical capacity using resource pools, allowing a workload to scale without necessarily adding a separate service instance. \\\addlinespace[2pt]
Service load balancing & The architecture specializes workload distribution for service endpoints and may use external or built-in balancing with a pool of equivalent service implementations. \\\addlinespace[2pt]
Cloud bursting & Monitoring detects local saturation, cloud capacity becomes active, traffic or work is redirected, and remote resources are released after the burst. \\\addlinespace[2pt]
Elastic disk provisioning & Thin or elastic provisioning exposes an apparent capacity while physical blocks are allocated on demand and usage is monitored. \\\addlinespace[2pt]
Redundant storage & A gateway or storage system coordinates copies and redirects access when a primary path or device fails, subject to consistency and recovery policy. \\\addlinespace[2pt]
Multicloud & Workloads, data, or supporting services are allocated across providers whose delivery models, interfaces, controls, regions, and operating practices can differ. \\\addlinespace[2pt]
Architecture selection & The practitioner identifies the problem, constraints, candidate pattern, participating mechanisms, expected improvement, tradeoffs, and operating evidence. \\\addlinespace[2pt]
\bottomrule
\end{longtable}
\begin{takeawaybox}
Cloud architecture benefits are realized only when the mechanisms, triggers, state, capacity, failure domains, and operating responsibilities behind the pattern are verified.
\end{takeawaybox}
\end{document}
